\documentclass[11pt,a4paper]{article}

\usepackage[T1]{fontenc}
\usepackage{lmodern}
\usepackage{amsmath,amssymb,amsthm}
\usepackage{booktabs,array}
\usepackage{enumitem}
\usepackage{xcolor}
\usepackage{pict2e}
\usepackage{xurl}
\usepackage[hidelinks]{hyperref}

\hypersetup{
  pdftitle={Color-split frame reductions for queen domination on even boards of sizes 50 through 64},
  pdfauthor={Edwin Chan}
}

\setlength{\textwidth}{\dimexpr\paperwidth-2.10in\relax}
\setlength{\textheight}{\dimexpr\paperheight-2.05in\relax}
\setlength{\oddsidemargin}{0.05in}
\setlength{\evensidemargin}{\oddsidemargin}
\setlength{\headheight}{12pt}
\setlength{\headsep}{16pt}
\setlength{\topmargin}{\dimexpr-0.06in-\headheight-\headsep\relax}
\setlength{\footskip}{26pt}
\setlength{\parindent}{1.1em}
\setlength{\parskip}{1.8pt plus 0.5pt}
\setlength{\emergencystretch}{1.3em}
\linespread{1.06}
\raggedbottom
\clubpenalty=10000
\widowpenalty=10000
\displaywidowpenalty=10000
\setlength{\jot}{5pt}
\AtBeginDocument{%
  \setlength{\abovedisplayskip}{11pt plus 2pt minus 2pt}%
  \setlength{\belowdisplayskip}{11pt plus 2pt minus 2pt}%
  \setlength{\abovedisplayshortskip}{6pt plus 2pt minus 1pt}%
  \setlength{\belowdisplayshortskip}{9pt plus 2pt minus 1pt}%
}

\makeatletter
\def\ps@preprint{%
  \let\@mkboth\@gobbletwo
  \def\@oddhead{\normalfont\small Queen domination numbers\hfil\thepage}%
  \let\@evenhead\@oddhead
  \let\@oddfoot\@empty
  \let\@evenfoot\@empty
}
\pagestyle{preprint}
\renewcommand{\@seccntformat}[1]{\csname the#1\endcsname.\enspace}
\renewcommand{\section}{\@startsection{section}{1}{\z@}%
  {-21pt plus -4pt minus -2pt}{8pt plus 2pt}%
  {\normalfont\fontsize{14}{17}\selectfont\bfseries\raggedright}}
\renewcommand{\subsection}{\@startsection{subsection}{2}{\z@}%
  {-15pt plus -3pt minus -2pt}{5pt plus 1pt}%
  {\normalfont\bfseries\raggedright}}
\renewenvironment{proof}[1][\proofname]{%
  \par\pushQED{\qed}\normalfont
  \topsep8pt plus 2pt minus 1pt\relax
  \trivlist\item[\hskip\labelsep\itshape #1\@addpunct{.}]\ignorespaces
}{\popQED\endtrivlist\@endpefalse}
\def\@maketitle{%
  \begingroup
    \raggedright
    {\fontsize{17}{21}\selectfont\bfseries\@title\par}%
    \vskip9pt
    {\normalsize\@author\par}%
    \vskip3pt
    {\small\@date\par}%
  \endgroup
  \vskip14pt
}
\renewenvironment{abstract}{%
  \par\noindent\textbf{Abstract}\par\nobreak\vskip3pt
  \noindent\ignorespaces
}{\par}
\setlength{\abovecaptionskip}{0pt}
\setlength{\belowcaptionskip}{8pt}
\long\def\@makecaption#1#2{%
  \vskip\abovecaptionskip
  {\small\raggedright\noindent\textbf{#1.}\enspace #2\par}%
  \vskip\belowcaptionskip
}
\makeatother

\newtheoremstyle{result}
  {10pt plus 2pt minus 1pt}{8pt plus 2pt minus 1pt}
  {\itshape}{}{\bfseries}{.}{0.55em}{}
\theoremstyle{result}
\newtheorem{theorem}{Theorem}[section]
\newtheorem{proposition}[theorem]{Proposition}
\newtheorem{lemma}[theorem]{Lemma}
\newtheorem{corollary}[theorem]{Corollary}
\theoremstyle{definition}
\newtheorem{definition}[theorem]{Definition}
\theoremstyle{remark}
\newtheorem{remark}[theorem]{Remark}

\newcommand{\waste}{\mathrm{waste}}
\newcommand{\Bd}{\mathrm{Bd}}
\newcommand{\Stab}{\mathrm{Stab}}
\newcommand{\ReproCommit}{61e8df4438cad8837f4498ad4160394028ca7232}
\newenvironment{placement}[1]{%
  \par\addvspace{7pt}\begin{samepage}%
  \noindent\textbf{#1}\par\nobreak\noindent\ignorespaces
}{\par\end{samepage}}

\title{Color-split frame reductions for queen domination\\
  on even boards of sizes 50 through 64}
\author{Edwin Chan\thanks{Gemini~4~Neon carried out the original research in Google Antigravity. See Section~\ref{sec:provenance} for contributions.}}
\date{November 27, 2026}

\begin{document}
\maketitle

\begin{abstract}
For every even $n$ from $50$ through $64$, the queen's graph $Q_n$ has
domination and independent domination numbers
$\gamma(Q_n)=i(Q_n)=n/2+1$. To prove the lower bounds, we group
$n/2$-queen placements by their outermost empty rows and columns and
refine the resulting counts by chessboard color. We refute all
$290\,542$ SAT instances with DRAT proofs checked by \texttt{drat-trim}.
The published independent placements of \"Osterg\aa rd and Weakley
give matching upper bounds for both parameters. We also give a
$31$-queen independent placement on $Q_{60}$ and report the time,
memory, and storage used by the computations.
\end{abstract}

\section{Introduction}\label{sec:intro}

The queen's graph $Q_n$ has the squares of an $n\times n$ chessboard as
vertices. Distinct squares are adjacent if they share a row, column, or
diagonal. The domination number $\gamma(Q_n)$ is the least size of a
set of queens that occupies or attacks every square. The independent
domination number $i(Q_n)$ is the minimum size of such a set with
pairwise non-attacking queens. Thus $\gamma(Q_n)\le i(Q_n)$.

For even $n$, the classical lower bound is $\gamma(Q_n)\ge n/2$,
due to Spencer (as reported by Cockayne~\cite{Cockayne1990}) and
Raghavan and Venkatesan~\cite{RV1987}. For even $26\le n\le120$,
the constructions of \"Osterg\aa rd and Weakley
\cite[Section~5]{OW2001} give $i(Q_n)\le n/2+1$.
We determine which of the two values occurs on the eight even
boards from $Q_{50}$ through $Q_{64}$.

\begin{theorem}\label{thm:main}
For every even $50\le n\le64$, $\gamma(Q_n)=n/2+1$. Explicitly,
\[
  \begin{aligned}
    \gamma(Q_{50})&=26, & \gamma(Q_{52})&=27, & \gamma(Q_{54})&=28, & \gamma(Q_{56})&=29,\\
    \gamma(Q_{58})&=30, & \gamma(Q_{60})&=31, & \gamma(Q_{62})&=32, & \gamma(Q_{64})&=33.
  \end{aligned}
\]
\end{theorem}

\begin{proof}
A dominating set of size at most $k=n/2$ extends to one of size $k$
by adding squares. By Theorem~\ref{thm:coverage}, some image of
this set under a board symmetry would extend to a satisfying
assignment of a generated SAT instance. The DRAT checks in
Section~\ref{sec:results} refute every instance on the eight boards.
The placements in Appendix~\ref{sec:witnesses} give the upper bounds.
\end{proof}

\begin{corollary}\label{cor:indep}
For every even $50\le n\le64$, $i(Q_n)=\gamma(Q_n)=n/2+1$.
\end{corollary}

\begin{proof}
Theorem~\ref{thm:main} gives $n/2+1=\gamma(Q_n)\le i(Q_n)$.
The published independent placements~\cite[Section~5, p.~11]{OW2001}
give the reverse bound. See also Weakley~\cite[Proposition~21]{Weakley2022}
and the coordinate checks in Appendix~\ref{sec:witnesses}.
\end{proof}

\newpage
\subsection{Previous work and proof outline}

The two problems go back to de~Jaenisch's 1863 appendix
\cite[2017 corrigendum]{OW2001} and appear as Problem~C18 in
Guy~\cite{Guy1994}. Their sequences in the OEIS are
\href{https://oeis.org/A075458}{A075458} and
\href{https://oeis.org/A075324}{A075324}, respectively.

Neuhaus~\cite[Section~2.4 and Chapter~4]{Neuhaus2009} reported
integer-programming results for several even boards through $n=46$
and for selected larger boards. His analysis of frame sides used
parity to track excess diagonal coverage. Weakley
\cite[Sections~2--3]{Weakley2022} obtained balance conditions for
bichromatic sets and exclusions for monochromatic ones. The
\texttt{BSC} Lean proof for $Q_{26}$ also uses parity and polynomial
identities~\cite{BSC2026}.

\emph{Hint and Conquer}~\cite{Hint2026} split the SAT search by
the outermost empty rows and columns and used side counts to
certify even values through $n=38$. Davies
\cite[Sections~2 and~4]{Davies2026} refined this split by the
deficits in the numbers of occupied lines and certified further
even values through $n=46$. His even-board runs used the default
CaDiCaL~3.1.0 configuration without symmetry breaking or cubing.
Table~2 of that paper reports single-core CPU times on a shared
machine. \DJ uki\'c~\cite{Djukic2026} reported a DRAT-certified
computation for $n=48$. On the even boards with $32\le n\le48$,
these ordinary values and the published independent placements
also give $i(Q_n)=n/2+1$.

Our reduction counts border and side incidences separately for
the two chessboard colors. These counts rule out some frames
algebraically and split the rest into SAT instances with tighter
constraints. Sections~\ref{sec:prelim}--\ref{sec:encoding} derive
the reductions and prove that the instances cover every possible
dominating $n/2$-set. Section~\ref{sec:results} reports the checks,
and Appendix~\ref{sec:witnesses} lists the placements.

\section{Frame geometry and dihedral symmetry}\label{sec:prelim}

Coordinates $(x,y)\in\{0,\ldots,n-1\}^2$ index columns and rows.
The difference and sum diagonals through $(x,y)$ have indices
$t=y-x$ and $s=x+y$, respectively. The chessboard color of a square
is $(x+y)\bmod 2$, taken in $\{0,1\}$. Each diagonal has the color
of its index modulo two, since $y-x\equiv x+y\pmod 2$.

Unless stated otherwise, take $n=2k\ge4$ and let $D$ be a
dominating set of $k$ queens. Write $R,C$ for the occupied rows and
columns and $T,S$ for the occupied difference and sum diagonals.
We identify each line by its type and index, so $T$ and $S$ are
disjoint. For each line family the \emph{deficit} is the number of
queens minus the number of occupied lines.
\[
  d_R = k - |R|, \quad d_C = k - |C|, \quad d_T = k - |T|, \quad d_S = k - |S| \qquad (\text{all } \ge 0).
\]
At least two rows and two columns are empty because
$n-|R|,n-|C|\ge k\ge2$. Let $a<b$ be the outermost empty columns
and $c<d$ the outermost empty rows (Figure~\ref{fig:frame}).
Use integer intervals throughout.
\begin{itemize}[leftmargin=1.8em,itemsep=2pt]
  \item The margins are $m_l=a$, $m_r=n-1-b$, $m_b=c$, and $m_t=n-1-d$.
  \item The frame has width $w=b-a\ge1$ and height $h=d-c\ge1$.
  \item The vertical sides are $V=\{a,b\}\times[c,d]$ and the horizontal
    sides are $H=[a,b]\times\{c,d\}$. The border $\Bd=V\cup H$ counts each
    corner once.
\end{itemize}

\begin{figure}[!htbp]
\centering
\setlength{\unitlength}{5.2mm}
\begin{picture}(14.75,9.8)(-1.15,-1.1)
  \begingroup
  \color{black!15}
  \multiput(1,0)(2,0){4}{\rule{1\unitlength}{1\unitlength}}
  \multiput(0,1)(2,0){4}{\rule{1\unitlength}{1\unitlength}}
  \multiput(1,2)(2,0){4}{\rule{1\unitlength}{1\unitlength}}
  \multiput(0,3)(2,0){4}{\rule{1\unitlength}{1\unitlength}}
  \multiput(1,4)(2,0){4}{\rule{1\unitlength}{1\unitlength}}
  \multiput(0,5)(2,0){4}{\rule{1\unitlength}{1\unitlength}}
  \multiput(1,6)(2,0){4}{\rule{1\unitlength}{1\unitlength}}
  \multiput(0,7)(2,0){4}{\rule{1\unitlength}{1\unitlength}}
  \endgroup
  \linethickness{0.2pt}
  \multiput(0,0)(1,0){9}{\line(0,1){8}}
  \multiput(0,0)(0,1){9}{\line(1,0){8}}
  \linethickness{0.95pt}
  \put(1.5,1.5){\framebox(5,4){}}
  \begingroup
  \color{blue!65!black}
  \put(1.5,1.5){\line(1,1){4}}
  \put(3.1,2.9){\makebox(0,0)[t]{\small $t$}}
  \endgroup
  \begingroup
  \color{red!60!black}
  \put(2.5,5.5){\line(1,-1){4}}
  \put(5.6,2.3){\makebox(0,0)[r]{\small $s$}}
  \endgroup
  \put(1.5,-0.27){\makebox(0,0)[t]{$a$}}
  \put(6.5,-0.27){\makebox(0,0)[t]{$b$}}
  \put(-0.24,1.5){\makebox(0,0)[r]{$c$}}
  \put(-0.24,5.5){\makebox(0,0)[r]{$d$}}
  \put(8.2,-0.31){\makebox(0,0)[l]{$x$}}
  \put(-0.21,8.23){\makebox(0,0)[r]{$y$}}
  \put(8.3,6.3){\makebox(0,0)[l]{\small Vertical sides $V$}}
  \put(8.3,5.5){\makebox(0,0)[l]{\small Horizontal sides $H$}}
  \put(8.3,4.1){\makebox(0,0)[l]{\small White squares ($q=0$)}}
  \put(8.3,3.3){\makebox(0,0)[l]{\small Shaded squares ($q=1$)}}
\end{picture}
\caption{The border $\Bd=V\cup H$ on an $8\times8$ board, drawn through
  the centers of the boundary squares. The difference diagonal $t=0$
  has color $0$ and the sum diagonal $s=7$ has color $1$.}\label{fig:frame}
\end{figure}

Every column outside $[a,b]$ and every row outside $[c,d]$ is occupied.
Thus
\begin{equation}\label{eq:margins}
  m_l + m_r \le |C| = k - d_C, \qquad m_b + m_t \le |R| = k - d_R.
\end{equation}
A square is \emph{needy} if its row and column are both empty.
Every needy square must lie on an occupied diagonal. The border
contains no queens because each side lies on an empty line.

Let $z_q$ count queens of color $q$, so $z_0+z_1=k$, and let
$T_q,S_q$ be the occupied diagonals of that color. If $d_T=0$
then $|T_q|=z_q$, and if $d_S=0$ then $|S_q|=z_q$. Write $e_p(R)$
and $e_p(C)$ for the numbers of empty rows and columns of parity $p$.
A placement using one color is \emph{monochromatic}. A placement
using both is \emph{bichromatic}.

The margins and deficits form the \emph{case} of $D$,
\[
  \kappa(D)=(m_l,m_r,m_b,m_t;\,d_R,d_C,d_T,d_S).
\]
Write $\rho(x)=n-1-x$. The following table gives the action of
the dihedral group $D_4$ on squares and on the case tuple.
\begin{align*}
  g_0 &: (x,y)\mapsto(x,y),           & (m_l,m_r,m_b,m_t;\,d_R,d_C,d_T,d_S),\\
  g_1 &: (x,y)\mapsto(\rho(x),y),     & (m_r,m_l,m_b,m_t;\,d_R,d_C,d_S,d_T),\\
  g_2 &: (x,y)\mapsto(x,\rho(y)),     & (m_l,m_r,m_t,m_b;\,d_R,d_C,d_S,d_T),\\
  g_3 &: (x,y)\mapsto(\rho(x),\rho(y)), & (m_r,m_l,m_t,m_b;\,d_R,d_C,d_T,d_S),\\
  g_4 &: (x,y)\mapsto(y,x),           & (m_b,m_t,m_l,m_r;\,d_C,d_R,d_T,d_S),\\
  g_5 &: (x,y)\mapsto(\rho(y),x),     & (m_t,m_b,m_l,m_r;\,d_C,d_R,d_S,d_T),\\
  g_6 &: (x,y)\mapsto(y,\rho(x)),     & (m_b,m_t,m_r,m_l;\,d_C,d_R,d_S,d_T),\\
  g_7 &: (x,y)\mapsto(\rho(y),\rho(x)), & (m_t,m_b,m_r,m_l;\,d_C,d_R,d_T,d_S).
\end{align*}
On an even board $g_1,g_2,g_5,g_6$ reverse chessboard parity and exchange
diagonal types. In a square frame they exchange the two corner-to-corner
diagonals. The other four maps preserve parity and diagonal types.
The maps $g_4,g_5,g_6,g_7$ exchange vertical and horizontal sides.

We take the lexicographically least of the eight case tuples.
This extends the frame-only canonicalization in
Davies~\cite[Sections~2 and~4]{Davies2026} to include the deficits.
Every placement has a dihedral image in a canonical case.
For such a case $\kappa$, the stabilizer
$\Stab(\kappa)=\{g\in D_4:g(\kappa)=\kappa\}$ gives the symmetries
still available within that case.

\section{Color-split frame and side identities}\label{sec:colorsplit}

We count how many needy squares the occupied diagonals cover on
the border and on each pair of sides. For $X\subseteq\Bd$, write
$X_q$ for its squares of color $q$ and $N_q(X)$ for the number
of needy squares in $X_q$. At a square $p$, let $m(p)\in\{0,1,2\}$
count occupied diagonals. For any diagonal $L$, let
$\beta_X(L)=|X\cap L|$.

Every diagonal meets $\Bd$ in at most two squares. A diagonal
meeting both vertical sides changes height by $w$, so when $w\ge h$
this can occur only at opposite corners in a square frame. The
horizontal statement follows by transpose.

The \emph{box diagonals} of a square frame are the difference
diagonal $U$ through $(a,c),(b,d)$, with $t=c-a$, and the sum
diagonal $W$ through $(a,d),(b,c)$, with $s=a+d$.
Write $[P]$ for the indicator of a proposition $P$. Set
\[
  u_0=[U\in T],\qquad v_0=[W\in S],\qquad
  \delta_q=[U\in T_q]+[W\in S_q].
\]
Set $u_0=v_0=\delta_q=0$ for nonsquare frames. The \emph{waste}
counts record unused coverage. They count missed intersections,
hits on squares already covered by a row or column, and second
hits on needy squares. The subscript $F$ denotes the full border.
\begin{align*}
  \waste_F^{(q)} &:= \sum_{L \in T_q \cup S_q} (2 - \beta_{\Bd}(L))\\
    &\quad + \#\{(L,p) : L \in T_q \cup S_q,\; p \in \Bd_q \cap L \text{ non-needy}\}\\
    &\quad + \#\{\text{needy } p \in \Bd_q : m(p)=2\}, \\
  \waste_V^{(q)} &:= \#\{L \in T_q \cup S_q : \beta_V(L)=0\}\\
    &\quad + \#\{(L,p) : L \in T_q \cup S_q,\; p \in V_q \cap L \text{ non-needy}\}\\
    &\quad + \#\{\text{needy } p \in V_q : m(p)=2\}.
\end{align*}
Define $\waste_H^{(q)}$ with $H$ in place of $V$. These counts
also make sense on odd boards whenever at least two rows and two
columns are empty.

\begin{theorem}[Border and side counts by color]\label{thm:colorsplit}
For every dominating set $D$ of $k$ queens on $Q_n$ with $k\le n-2$,
and each $q\in\{0,1\}$, we have
\begin{align}
  \waste_F^{(q)} &= 2(|T_q| + |S_q|) - N_q(\Bd) \ge 0, \label{eq:Fq} \\
  \waste_V^{(q)} &= (|T_q| + |S_q|) + \delta_q - N_q(V) \ge 0 \qquad (w \ge h), \label{eq:Vq} \\
  \waste_H^{(q)} &= (|T_q| + |S_q|) + \delta_q - N_q(H) \ge 0 \qquad (h \ge w). \label{eq:Hq}
\end{align}
\end{theorem}

\begin{proof}
Fix a color $q$. Let $I_q(X)$ count incidences between occupied
diagonals of color $q$ and squares in $X_q$. Each needy square
contributes at least one incidence, giving
\[
  \begin{aligned}
    I_q(X) = \sum_{p \in X_q} m(p)
      &= N_q(X) + \#\{\text{needy } p \in X_q : m(p)=2\}\\
      &\quad + \#\{(L,p) : L \in T_q \cup S_q,\; p \in X_q \cap L \text{ non-needy}\}.
  \end{aligned}
\]
Counting by diagonals gives $I_q(X)=\sum_{L\in T_q\cup S_q}\beta_X(L)$.
On the full border, this is twice the number of occupied diagonals
of color $q$ minus the first term of $\waste_F^{(q)}$. On $V$ when
$w\ge h$, each occupied diagonal contributes one hit, with zero
hits for those that miss $V$ and an extra hit for each occupied box
diagonal of color $q$. The corresponding count on $H$ applies when
$h\ge w$. Substituting these counts gives \eqref{eq:Fq}--\eqref{eq:Hq}.
Every term in the waste counts is nonnegative.
\end{proof}

Each empty row contributes two needy squares to $V$ and each empty
column two to $H$. All four corners are needy and counted in both,
so the totals over both colors are
\begin{align*}
  \sum_q N_q(V) &= 2(n-|R|), & \sum_q N_q(H) &= 2(n-|C|),\\
  \sum_q N_q(\Bd) &= 2(n-|R|)+2(n-|C|)-4.
\end{align*}
Write $\waste_F,\waste_V,\waste_H$ for the sums over both colors.
For $n=2k$ the identities reduce to those of
Davies~\cite[Section~2]{Davies2026}.
\begin{align}
  \waste_F &= 4 - 2(d_R + d_C + d_T + d_S) \ge 0, \label{eq:Ftot} \\
  \waste_V &= -d_T - d_S - 2d_R + [w=h](u_0+v_0) \ge 0 \qquad (w \ge h), \label{eq:Vtot} \\
  \waste_H &= -d_T - d_S - 2d_C + [w=h](u_0+v_0) \ge 0 \qquad (h \ge w). \label{eq:Htot}
\end{align}
In particular, $d_R + d_C + d_T + d_S \le 2$, so at least two
deficits are zero. The four are at most $k-1$, since a nonempty
placement occupies a line of each type.

\section{Algebraic exclusion of cases}\label{sec:algebraic}

The bounds on waste already restrict the deficits. Four further
tests exclude cases before the SAT formulas are built. The margin
test takes $O(k)$ steps per case, and each parity test takes a
constant number of steps.

\begin{lemma}[NS-1: Margin bounds for nonsquare frames]\label{lem:ns1}
Suppose $w>h$. Then $d_R=d_T=d_S=0$ and $\waste_V=0$, and
the margins satisfy the following bounds.
\begin{enumerate}[label=(\roman*),leftmargin=2em,itemsep=2pt]
  \item $m_l \le \lfloor h/2 \rfloor$ and $m_r \le \lfloor h/2 \rfloor$;
  \item $m_l + m_r + m_b + m_t \le k$;
  \item for every $r \in \{1, \dots, \max(m_l, m_r)\}$, $\max(0, m_l - r + 1) + \max(0, m_r - r + 1) \le h + 1 - 2r$;
  \item For each $m\in\{m_b,m_t\}$ with $m>0$, one has
    $2m\le w$ and $w\le 2h+2$.
\end{enumerate}
For $h>w$, transpose by $g_4$.
\end{lemma}

\begin{proof}
By \eqref{eq:Vtot}, $-d_T-d_S-2d_R\ge0$. Thus all three
deficits vanish and $\waste_V=0$. Every occupied diagonal meets
$V$, so each queen $(x,y)\in D$ satisfies
\[
  \begin{aligned}
    y-x &\in [c-b,d-b]\cup[c-a,d-a],\\
    x+y &\in [a+c,a+d]\cup[b+c,b+d].
  \end{aligned}
\]
For a margin column $x=a-\delta$, $\delta\ge1$, the first difference
interval gives $y\le d-w-\delta\le c-2$. Both sum intervals give
$y\ge c+\delta$. The second sum interval gives $y\ge c+w+\delta\ge d+2$,
while both difference intervals give $y\le d-\delta$.
The remaining intersection is
\[
  y \in [c + \delta,\; d - \delta] \subseteq (c, d),
\]
which requires $2\delta\le h$. On any side with a positive
margin, its outermost occupied column gives (i).

Choose one queen in each margin column. Since $d_R=0$, they
occupy distinct rows, all inside $(c,d)$. The $m_b+m_t$ occupied
rows outside $[c,d]$ bring the total to at least
$m_l+m_r+m_b+m_t$, giving (ii).

For (iii), the chosen queens in columns at distance at least $r$
from $[a,b]$ must fit in the $h+1-2r$ rows of $[c+r,d-r]$.
There are $\max(0,m_l-r+1)$ such columns on the left and
$\max(0,m_r-r+1)$ on the right.

For (iv), take a margin row $y=c-\delta$ with $\delta\ge1$.
The difference and sum intervals restrict $x$ to
$(X_{t,1}\cup X_{t,2})\cap(X_{s,1}\cup X_{s,2})$, where
\begin{align*}
  X_{t,1} &= [b - h - \delta,\; b - \delta], & X_{t,2} &= [a - h - \delta,\; a - \delta], \\
  X_{s,1} &= [a + \delta,\; a + h + \delta], & X_{s,2} &= [b + \delta,\; b + h + \delta].
\end{align*}
The interval $X_{t,2}$ lies strictly to the left of both sum
intervals, and $X_{t,1}$ ends before $X_{s,2}$ begins. Thus
\[
  x \in [\max(b - h - \delta,\; a + \delta),\; \min(b - \delta,\; a + h + \delta)],
\]
which is nonempty exactly when $2\delta\le w$ and $w-2h\le2\delta$.
If $m_b>0$, row $0$ gives the first bound with $\delta=m_b$,
and row $c-1$ gives the second with $\delta=1$. Reflection gives
the bounds for $m_t$.
\end{proof}

\begin{lemma}[NS-2: Nonsquare frames with an odd longer side]\label{lem:ns2}
If $k$ is odd, the longer side of a nonsquare frame has even length.
\end{lemma}

\begin{proof}
Transpose if needed and suppose $w>h$ with $w$ odd. The columns
$a,b$ have opposite parity, so each of the $k$ empty rows
contributes one needy square of each color to $V$.
Lemma~\ref{lem:ns1} gives $|T_q|=|S_q|=z_q$ and
$\waste_V^{(q)}=0$. Equation~\eqref{eq:Vq} then gives
$2z_q=k$, impossible for odd $k$.
\end{proof}

\begin{lemma}[NS-3: Mixed side parity and deficit two]\label{lem:ns3}
For odd $k$, a nonsquare frame with $w>h$, $w$ even, and $h$ odd
cannot have $d_C=2$. The transposed case with $d_R=2$ is also impossible.
\end{lemma}

\begin{proof}
Suppose $d_C=2$ in the stated orientation.
Lemma~\ref{lem:ns1} and \eqref{eq:Vq} give $N_q(V)=2z_q$.
Since $h$ is odd, each of the $k+2$ empty columns contributes
one needy square of each color to $H$. Two corners of each color
lie in both side sets, so $N_q(\Bd)=2z_q+k$. By \eqref{eq:Fq},
\[
  0 \le \waste_F^{(q)} = 4 z_q - (2 z_q + k) = 2 z_q - k \implies 2 z_q \ge k \qquad (q \in \{0,1\}).
\]
Since $k$ is odd, $z_0 \ge (k+1)/2$ and $z_1 \ge (k+1)/2$, contradicting $z_0 + z_1 = k$.
\end{proof}

\begin{lemma}[SQ-RC-Odd: A single row or column deficit]\label{lem:sqrc}
For odd $k$, a square frame of odd side length cannot have
deficit vector $(0,1,0,0)$ or $(1,0,0,0)$.
\end{lemma}

\begin{proof}
Transpose if needed and assume $d_C=1$. By \eqref{eq:Htot},
$\waste_H=-2+u_0+v_0\ge0$, so both box diagonals are occupied and
$\waste_H=0$. The odd side length puts one box diagonal in each
color and one needy square of each color on $H$ in every empty
column. Thus $N_q(H)=k+1$ and $\delta_q=1$. With $d_T=d_S=0$,
\eqref{eq:Hq} gives $2z_q+1=k+1$, again impossible for odd $k$.
\end{proof}

\section{Branches by parity and deficit}\label{sec:branches}

The remaining cases are either nonsquare or square with total
deficit at most two. Within each case we branch on the color
counts, the box diagonals, or both. Each resulting SAT instance
is called a \emph{leaf}.

Write $p_a=a\bmod 2$ and $p_c=c\bmod 2$, and use $\oplus$ for
addition modulo two. If $w,h$ are both even, all four corners
have color $q_*=p_a\oplus p_c$. Write $\bar q_*=1-q_*$.
For an odd-sided square, let $\eta_q$ indicate whether its box
diagonal of color $q$ is occupied.

\subsection{Nonsquare cases}
Transpose if needed so $w>h$. Then $d_R=d_T=d_S=0$,
$d_C\in\{0,1,2\}$, and $\waste_V^{(q)}=0$ for both colors.
\begin{itemize}[leftmargin=1.8em,itemsep=3pt]
  \item \textbf{Odd $w$.} By Lemma~\ref{lem:ns2}, $k$ is even.
    The vertical side equation gives $2z_q=N_q(V)=k$,
    so $|T_q|=|S_q|=k/2$ for each color. The border waste
    is at most $4-2d_C$ and vanishes when $d_C=2$.
  \item \textbf{Even $w$, odd $h$.}
    Here $N_q(V)=2e_{p_a\oplus q}(R)=2z_q$. Each empty column
    adds one square of color $q$ to $H$, with two corners already
    counted on $V$. Thus $N_q(\Bd)=2z_q+k+d_C-2$ and
  \[
    \waste_F^{(q)} = 2 z_q - k - d_C + 2 \ge 0 \qquad (q \in \{0,1\}).
  \]
    Branch on
    $z_0\in[\lceil(k+d_C-2)/2\rceil,\lfloor(k-d_C+2)/2\rfloor]$.
    There are at most two choices for odd $k$ and three for even
    $k$. If a case symmetry exchanges colors, keep $z_0\le k-z_0$.
    Each choice fixes the common count $|T_q|=|S_q|=e_{p_a\oplus q}(R)=z_q$,
    as well as the border waste of each color.
  \item \textbf{Even $w$ and $h$.}
    Put $z=z_{\bar q_*}$. Then $N_{\bar q_*}(V)=2e_{1-p_c}(R)=2z$.
    All corners have color $q_*$, so
    $N_{\bar q_*}(\Bd)=2e_{1-p_c}(R)+2e_{1-p_a}(C)$. With
    $\Delta=z-e_{1-p_a}(C)$ the two border wastes are
    \[
      \waste_F^{(\bar q_*)}=2\Delta,\qquad
      \waste_F^{(q_*)}=4-2d_C-2\Delta.
    \]
    Thus $\Delta\in\{0,\ldots,2-d_C\}$. The first waste vanishes
    at the lower endpoint and the second at the upper endpoint.
\end{itemize}

\subsection{Square frames with total deficit two}
The side inequalities leave only $(0,0,2,0)$, $(0,0,0,2)$,
$(0,0,1,1)$, and $(1,1,0,0)$, subject to the individual bound $k-1$.
In each case \eqref{eq:Ftot}--\eqref{eq:Htot} force
$u_0=v_0=1$ and zero border and side waste in both colors.

\subsection{Square frames with one row or column deficit}
Take $d_C=1$ and the other deficits zero. The case $d_R=1$
follows by transposition. Equation~\eqref{eq:Htot} gives
$u_0=v_0=1$ and $\waste_H=0$.
\begin{itemize}[leftmargin=1.8em,itemsep=3pt]
  \item \textbf{Odd $w$.} The argument of Lemma~\ref{lem:sqrc} gives
    $2z_q=k$, so $k$ is even. The vertical and border equations
    give $\waste_V^{(q)}=\waste_F^{(q)}=1$ for each color.
  \item \textbf{Even $w$.} Zero horizontal waste in color $\bar q_*$
    gives $z_{\bar q_*}=e_{1-p_a}(C)$. Put
    $\Delta=z_{\bar q_*}-e_{1-p_c}(R)$. Equations
    \eqref{eq:Vq} and~\eqref{eq:Fq} give
    \[
      \waste_V^{(\bar q_*)}=\waste_F^{(\bar q_*)}=2\Delta,\qquad
      \waste_V^{(q_*)}=\waste_F^{(q_*)}=2-2\Delta.
    \]
    Hence $\Delta$ is zero or one. Color $\bar q_*$ has zero
    waste on all three sets when $\Delta=0$, and color $q_*$
    has zero waste when $\Delta=1$.
\end{itemize}

\subsection{Square frames with one diagonal deficit}
Take $d_S=1$ and the other deficits zero. For $d_T=1$, exchange
$T$ and $S$. Every queen occupies a distinct difference diagonal,
so $|T_q|=z_q$. The one sum diagonal occupied twice has color
$q_{\mathrm{def}}$, and $|S_q|=z_q-[q=q_{\mathrm{def}}]$.
We branch on $q_{\mathrm{def}}$.
\begin{itemize}[leftmargin=1.8em,itemsep=3pt]
  \item \textbf{Odd $w$.} Let $q'=1-q_{\mathrm{def}}$. By \eqref{eq:Vq},
    \[
      2z_{q_{\mathrm{def}}}-1+\eta_{q_{\mathrm{def}}}\ge k,\qquad
      2z_{q'}+\eta_{q'}\ge k .
    \]
    For odd $k$, the inequalities force
    $z_{q_{\mathrm{def}}}=(k+1)/2$, $z_{q'}=(k-1)/2$, and $\eta_{q'}=1$.
    The color $q'$ has zero waste on the border and both sides.
    For even $k$, they force $z_0=z_1=k/2$ and $\eta_{q_{\mathrm{def}}}=1$.
    In this case $q_{\mathrm{def}}$ has zero waste on all three sets.
    In each case the other color has border waste two and side
    wastes equal to its box-diagonal indicator.
  \item \textbf{Even $w$.} Both box diagonals have color $q_*$. Put $z=z_{\bar q_*}$,
  $e_V=e_{1-p_c}(R)$, $e_H=e_{1-p_a}(C)$, and
  $\epsilon=[q_{\mathrm{def}}=\bar q_*]$. No corner has color $\bar q_*$,
  so \eqref{eq:Fq}--\eqref{eq:Hq} give
  \begin{align*}
    \waste_V^{(\bar q_*)}&=2z-\epsilon-2e_V,\\
    \waste_H^{(\bar q_*)}&=2z-\epsilon-2e_H,\\
    \waste_F^{(\bar q_*)}&=\waste_V^{(\bar q_*)}+\waste_H^{(\bar q_*)}.
  \end{align*}
  The two side wastes are at most $u_0+v_0-1\le1$.
  If $\epsilon=0$, both are even and hence zero, as is the border waste
  of color $\bar q_*$. Thus $e_V=e_H=z$, and $u_0+v_0\ge1$.
  The color $q_*$ then has border waste $2$ and side wastes
  $u_0+v_0-1$.
  If $\epsilon=1$, both are odd and hence one. This forces
  $u_0=v_0=1$, $e_V=e_H=z-1$, and zero waste of color $q_*$ on all
  three sets. For $d_S=1$ the signed count difference is
  $|T_{\bar q_*}|-1=|S_{\bar q_*}|=e_V=e_H$.
  When $d_T=1$, the roles of $T$ and $S$ are exchanged.
\end{itemize}

\subsection{Square frames with zero deficits}\label{sec:sq0}

\begin{proposition}[Branch parameters for square frames]\label{prop:sq0}
Suppose $w=h$ and all four deficits are zero.
\begin{enumerate}[label=(\alph*),leftmargin=2em,itemsep=3pt]
  \item \textbf{Odd $w$ and $k$.}
    The color counts are $(k-1)/2$ and $(k+1)/2$. The smaller
    color has its box diagonal occupied and zero waste on the
    border and both sides. If a case symmetry exchanges colors,
    choose color zero as the smaller color.
  \item \textbf{Odd $w$, even $k$.}
    Both color counts are $k/2$. For each color $q$,
    $\waste_V^{(q)}=\waste_H^{(q)}=\eta_q$ and
    $\waste_F^{(q)}=2$. When $\eta_q=1$, the border waste is
    the sum of the two side wastes. Branch on $(u_0,v_0)$,
    dropping $(0,1)$ if a case symmetry exchanges the two bits.
  \item \textbf{Even $w$.}
    Let $\Delta_V=z_{\bar q_*}-e_{1-p_c}(R)$ and
    $\Delta_H=z_{\bar q_*}-e_{1-p_a}(C)$. Both lie in $\{0,1\}$,
    with $u_0+v_0\ge2\max(\Delta_V,\Delta_H)$. This gives seven
    parameter tuples $(\Delta_V,\Delta_H,u_0,v_0)$ before the
    monochromatic refinement:
  \[
    (0,0,0,0),\; (0,0,1,0),\; (0,0,0,1),\; (0,0,1,1),\; (1,0,1,1),\; (0,1,1,1),\; (1,1,1,1).
  \]
    If $g_4$ or $g_7$ stabilizes the case, keep $(1,0,1,1)$
    as the representative of $(0,1,1,1)$. Each side has
    zero waste in at least one color. The border has zero
    waste in color $\bar q_*$ if both offsets are zero,
    and in color $q_*$ if both are one.
\end{enumerate}
\end{proposition}

\begin{proof}
For odd $w$, each empty row contributes one square of each color
to $V$, and each empty column does the same on $H$. The two corners
of each color are shared, giving $N_q(V)=N_q(H)=k$ and
$N_q(\Bd)=2k-2$. Thus the side inequalities become
$2z_q+\eta_q\ge k$. When $k$ is odd, they force the counts in (a)
and $\eta_q=1$ for the smaller color. Substituting into
\eqref{eq:Fq}--\eqref{eq:Hq} gives zero wastes for the smaller
color. For the larger color it gives border waste $4$ and side
wastes $1+\eta_q$. When $k$ is even the inequalities force the
equal counts in (b), so the side wastes equal $\eta_q$ and the
border waste is $2$.

For even $w$, the box diagonals have color $q_*$. The vertical
side contains $2e_{1-p_c}(R)$ needy squares of color $\bar q_*$.
The two side equations and the total waste give
\[
  \waste_X^{(\bar q_*)}=2\Delta_X,\qquad
  \waste_X^{(q_*)}=u_0+v_0-2\Delta_X,\qquad X\in\{V,H\}.
\]
Both wastes are nonnegative, giving the stated bounds on the
offsets and the zero wastes on the sides. The border wastes are
$2(\Delta_V+\Delta_H)$ in color $\bar q_*$ and
$4-2(\Delta_V+\Delta_H)$ in color $q_*$.
The symmetry actions in Section~\ref{sec:prelim} exchange colors
and box-diagonal indicators in (a) and (b), respectively.
In (c), $g_4$ and $g_7$ exchange the offsets and fix the box diagonals.
\end{proof}

\subsection{Monochromatic frames and diagonal moments}\label{sec:mono}

Weakley~\cite[Proposition~13]{Weakley2022} gives the row and column
structure of monochromatic dominating sets. The following lemma
identifies the corresponding frame in the case enumeration.

\begin{lemma}[Uniqueness of the monochromatic frame]\label{lem:mono}
Suppose $w=h$ is even, $d_R=d_C=0$, $d_T+d_S\le1$, and
$z_{\bar q_*}=0$. In the canonical frame, the margins are
$(0,1,0,1)$, the empty rows and columns have indices
$0,2,\ldots,n-2$, and the queens are at $(2i+1,2j+1)$ with
$0\le i,j<k$ and $j=\pi(i)$ for a permutation $\pi$.
\end{lemma}

\begin{proof}
For zero diagonal deficits, the offsets in
Proposition~\ref{prop:sq0}(c) are nonnegative. With one diagonal
deficit, the repeated diagonal has color $q_*$ and the
opposite-color empty-line counts equal $z_{\bar q_*}$. In either
case $z_{\bar q_*}=0$ forces $e_{1-p_c}(R)=e_{1-p_a}(C)=0$.
The $k$ empty lines of each family are therefore the entire even
or odd parity class. Each margin pair is $(0,1)$ or $(1,0)$.
Reflecting the board gives the canonical tuple $(0,1,0,1)$, so
$R=C=\{1,3,\ldots,2k-1\}$. The zero row and column deficits give
the permutation.
\end{proof}

On this $k\times k$ grid, set $t'=j-i$ and $s'=i+j+1$. Write
$T',S'$ for the occupied indices in $[-(k-1),k-1]$ and $[1,2k-1]$,
respectively.

The moment identities of \"Osterg\aa rd and Weakley
\cite[Theorem~4 and Eq.~(3)]{OW2001}, also used by Weakley
\cite[Eqs.~(1)--(2)]{Weakley2022}, take the following form here.
\begin{align*}
  d_T = 0 &\implies \sum_{t' \in T'} t' = 0, \qquad
  d_S = 0 \implies \sum_{s' \in S'} s' = k^2, \\
  d_T = d_S = 0 &\implies \sum_{t' \in T'} (t')^2 + \sum_{s' \in S'} (s' - 1)^2 = 4 \sum_{i=0}^{k-1} i^2 = \frac{2}{3}k(k-1)(2k-1).
\end{align*}
The two linear identities follow by summing $j-i$ and $i+j+1$
over the permutation. With both diagonal deficits zero, each summand
occurs once in its occupied family, and
$(j-i)^2+(i+j)^2=2(i^2+j^2)$ gives the quadratic identity. With one
diagonal deficit, we use the linear identity for the family with
deficit zero. The implementation imposes the available identities
modulo small integers; see Section~\ref{sec:encoding}.

\begin{remark}
The incidence identities of Theorem~\ref{thm:colorsplit} apply to odd
boards as well. The subsequent branch enumeration and the computations
reported here use $n=2k$.
\end{remark}

\section{SAT encoding and completeness}\label{sec:encoding}

For each branch we build a Boolean formula in conjunctive normal
form (CNF). Its clauses require a placement that dominates the
board and satisfies the branch constraints. We first bound the size
of the counters used in those constraints.

\begin{lemma}[Bounds for the parity counters]\label{lem:caps}
Write $z=z_{\bar q_*}$ in an even-by-even frame. The parity
equalities of Section~\ref{sec:branches} admit the following bounds.
\begin{enumerate}[label=(\alph*),leftmargin=2em,itemsep=2pt]
  \item For a nonsquare frame, $z\le k-2$.
  \item For a square with one row or column deficit, $z\le k-1$.
  \item For a square with one diagonal deficit, $z\le k-2$ when
    $q_{\mathrm{def}}=q_*$, and $z-1\le k-2$ otherwise.
  \item For a square with zero deficits,
    $z\le k-2+\max(\Delta_V,\Delta_H)$.
\end{enumerate}
\end{lemma}

\begin{proof}
A side of even length $\ell$ has exactly $\ell/2$ interior line
indices of parity opposite to its endpoints. In (a) the shorter
side is at most $2k-4$, and the number of empty lines of that parity
on the shorter side equals $z$. In (b), the corresponding equality
on the side with zero waste in both colors gives a bound of
$w/2\le k-1$. In (c), a repeated diagonal of color $q_*$ uses two
queens of that color. If the repeat is of the other color, both
box diagonals are occupied, so at least one queen has color $q_*$.

For (d), a needy corner requires at least one queen of color $q_*$.
This gives $z\le k-1$ and suffices when either $\Delta$ is positive.
When both vanish and $w\le 2k-4$, the side count suffices. The
remaining frame has $w=2k-2$ and, up to symmetry, margins
$(0,1,0,1)$. If $z=k-1$, all $k-1$ odd interior rows and columns
are empty, leaving a single available odd index $2k-1$ in each
direction. The queens in color $\bar q_*$ use exactly one odd
coordinate each. At most two fit in distinct rows and columns.
For $k\ge4$ this contradicts $z=k-1$. For $k=2$ both even indices
are empty boundaries, making any such queen impossible. For $k=3$,
the only available even index is $2$ in either direction. The two
queens would be $(5,2)$ and $(2,5)$ on the same sum diagonal,
contrary to zero deficit.
\end{proof}

\subsection{Board and auxiliary variables}

The Boolean variable $Q_{xy}$ records whether a queen occupies an
available square $(x,y)$. Squares on the four empty boundary lines
or on diagonals excluded by a zero-waste constraint are unavailable.
The variables $o_R(y)$, $o_C(x)$, $o_T(t)$, and $o_S(s)$ record
whether a row, column, difference diagonal, or sum diagonal contains
a queen. Each is the disjunction of its square variables, and is
false when the line has no available square.

For each square in the closed frame, a coverage clause requires
at least one of its four line variables to be true. Unit clauses
force the exterior rows and columns to be occupied, covering all
squares outside the frame.

The four occupancy counts are $k-d_R,k-d_C,k-d_T,k-d_S$. When
the diagonal counts are specified by color, the two counts add to
the same total. For each family with deficit zero, the encoder
allows at most one queen per line. Such a family always exists,
and its $k$ occupied lines force exactly $k$ queens.
The branch fixes counts of empty lines by parity, imposes
the relevant box-diagonal occupancies, and bounds the wastes.
In an even-sided zero-deficit square, if a box diagonal is empty,
the other diagonal at each of its corners must cover that corner.

The parity counters use the caps in Lemma~\ref{lem:caps}, tightened
by the lengths and offsets of the corresponding lists.

At a border square $p$, let $N_p$ indicate that the square is
non-needy. It is the row occupancy on a vertical side, the column
occupancy on a horizontal side, and zero at a corner. Write $o_L$
for the occupancy of diagonal $L$. The waste constraints count
the following literals, with weights where indicated.
\begin{itemize}[leftmargin=1.8em,itemsep=2pt]
  \item Missed intersections contribute $o_L$ with weight
    $2-\beta_{\Bd}(L)$ for the border or $[\beta_X(L)=0]$
    for a side set $X$ covered by Theorem~\ref{thm:colorsplit}.
  \item A hit on a non-needy square contributes $E_{L,p}$,
    forced true by $o_L\land N_p$.
  \item A double hit on a needy square contributes $F_p$, forced
    true by $o_T\land o_S\land\neg N_p$ at $p$.
\end{itemize}
A placement extends by assigning $E_{L,p}$ and $F_p$ their
conjunction values. For a zero-waste bound the encoder forbids
each corresponding event directly. It omits those same per-square
events from the border count if they were already forbidden on a
side. To encode a bound such as $\waste_V^{(q)}\le u_0+v_0-1$,
it adds $1-u_0$ and $1-v_0$ to the count and bounds the sum by one.

Cardinality constraints use unary totalizers~\cite{BB2003}.
A bit at index $i$ marks the threshold of $i$ true literals
in a group. When merging two groups, the forward clauses
propagate true child thresholds $i$ and $j$ to threshold $i+j$
of the combined count. A zero threshold is always true. At an
upper bound $B$, clauses forbid the first overflow $i+j=B+1$.
Inductively any count exceeding $B$ violates a clause.

When two counts must differ by a specified offset, reverse
clauses also make the parent thresholds false above the true count.
The two threshold vectors can then be equated after shifting by
the offset. For an exact constant count, an upper bound on the
negated literals supplies the matching lower bound. At-most-one
constraints along a line use pairwise exclusion or a sequential
counter~\cite{Sinz2005}, whose prefix bit indicates whether an
earlier prefix contains a queen. Assigning bits from the actual
counts satisfies all these clauses.

\subsection{Monochromatic formulas}

For the frame in Lemma~\ref{lem:mono}, the reduced encoder
places one queen in every row and column of the odd-odd
$k\times k$ grid. It indexes diagonal occupancy on color-zero
squares by $t'=t/2$ and $s'=s/2$.

The only needy squares are $(2u,2v)$ with $0\le u,v<k$.
At each of them the coverage clause is $o_T(v-u)\lor o_S(u+v)$.
Since $o_S(0)=0$, the corner $(0,0)$ forces $o_T(0)=u_0=1$.

With one diagonal deficit, \eqref{eq:Ftot}--\eqref{eq:Htot}
give border waste $2$ and side waste $v_0=o_S(k-1)$ on
each side set. The diagonals $|t'|=k-1$ and $s'=2k-2$
meet the frame border once, and $s'=2k-1$ misses it.
All other allocated diagonals meet the border twice and
hit both side sets. All border squares of the queen color
are needy, so the reduced waste bounds use only occupancy
and double-coverage bits.

The encoder imposes the moment identities modulo $3,4,5,7,8$
and, for zero deficit, also modulo $9$. At each partial sum it
uses one bit per possible residue. The clauses start the residue
at zero, add the next term's weight if its literal is true, and
require the final residue to equal the target. A permutation
satisfying the identities gives the required residues by direct
summation. The transpose sends $t'$ to $-t'$ and fixes $s'$,
allowing a lexicographic comparison of the difference-diagonal
vectors.

\subsection{Coverage of the generated formulas}

\begin{theorem}\label{thm:coverage}
For $n=2k\ge4$, every dominating $k$-set has a dihedral image whose
queen and occupancy indicators extend to a satisfying assignment
of a generated formula. Conversely, every satisfying assignment of
a generated formula gives a dominating set of size $k$ on $Q_n$.
\end{theorem}

\begin{proof}
Take a minimum case of the placement under the eight maps of
Section~\ref{sec:prelim}. The generator enumerates nonnegative
margin tuples satisfying \eqref{eq:margins} and $w,h\ge1$, and
deficit tuples with each entry at most $k-1$ and sum at most two.
Its side tests are the inequalities \eqref{eq:Vtot} and
\eqref{eq:Htot} with $u_0+v_0$ replaced by its upper bound $2$.
Hence the minimum case is enumerated. Lemmas~\ref{lem:ns1}--\ref{lem:sqrc}
show that no algebraic pruning rule discards it.

The actual parity counts and occupancy bits satisfy one of the
branch descriptions. A color-flipping case stabilizer exchanges
$z_0$ and $z_1$, exchanges the smaller color on an odd-sided
zero-deficit square when $k$ is odd, and exchanges the two box
diagonals there when $k$ is even. These are precisely the
identifications used to omit a representative in those
branches. Transposition and antitransposition interchange
$\Delta_V,\Delta_H$ in an even-sided square and fix the box
diagonals individually, justifying the omission of $(0,1,1,1)$
when one of those maps stabilizes the case.

For the two monochromatic patterns the conditions of
Lemma~\ref{lem:mono} force the special frame $(0,1,0,1)$.
The reduced encoder covers all such placements. The general
branch at that frame for a diagonal deficit in $q_*$ and the
zero-deficit branch $(0,0,1,1)$ both require $z_{\bar q_*}>0$,
since their zero-count placements are covered by the reduced
formula. Other zero-deficit branches may overlap with it.

Within a general formula the encoder compares the combined row
and column occupancy vector with its images under selected
members of $\Stab(\kappa)$. Color-reversing maps are kept only
when their action preserves the branch parameters. In the even-sided
square, an exchange of sides is kept only when it preserves the
fixed deficits and the ordered pair $(\Delta_V,\Delta_H)$ when
present. This follows from the eight actions on lines above. The
single-row-or-column-deficit even square uses only the half-turn.
In each comparison a prefix bit is forced when the vectors agree
up to that entry. The bit then activates the inequality at the
next entry. A placement with lexicographically least occupancy
vector in the orbit under the generated subgroup satisfies every
comparison. Its prefix bits can be set from the two vectors. The
reduced formula uses the same argument with transpose.

The resulting placement satisfies the coverage and cardinality
constraints. The branch restrictions follow from the incidence
identities, Lemma~\ref{lem:caps}, and the monochromatic moment
identities where applicable. Its actual occupancy, waste counts,
counter thresholds, and residues provide the auxiliary assignment.
Conversely, a satisfying assignment defines a dominating set by
the occupancy and coverage clauses, with exactly $k$ queens by
the zero-deficit family or the reduced permutation constraints.
\end{proof}

\section{Computations and verification}\label{sec:results}

\subsection{Certified runs and audit records}

All $290\,542$ formulas generated for the eight boards were
unsatisfiable. For each formula, the worker wrote a DIMACS file,
computed its SHA-256 digest, and requested a binary DRAT proof
from CaDiCaL~\cite{CaDiCaL}. The unmodified \texttt{drat-trim}
checker~\cite{DRATtrim} accepted every proof in backward mode.
The counts and timings are given in Table~\ref{tab:certified}.

\begin{table}[!htbp]
\centering\small
\caption{Certified runs on one ten-core Apple~M4 with ten workers
  and \texttt{-{}-config auto}. Solve and check times sum elapsed
  seconds across workers. Wall time measures the worker phase.
  Timing totals sum the reported board times. Proof volumes are in
  decimal gigabytes.}
\label{tab:certified}
\setlength{\tabcolsep}{4.2pt}
\renewcommand{\arraystretch}{1.18}
\begin{tabular}{@{}rrrrrrrrr@{}}
\toprule
$n$ & $k$ & Canonical & Pruned (\%) & Leaves & Solve (s) & Check (s) & Wall (s) & DRAT (GB) \\
\midrule
50 & 25 & 54\,626 & 44\,288 (81.1) & 19\,090 & 1\,799.29 & 2\,703.82 & 492.68 & 2.69 \\
52 & 26 & 63\,119 & 43\,354 (68.7) & 30\,814 & 5\,539.94 & 5\,434.66 & 1\,170.76 & 6.21 \\
54 & 27 & 72\,268 & 59\,164 (81.9) & 24\,251 & 3\,512.24 & 4\,388.23 & 856.67 & 4.35 \\
56 & 28 & 82\,695 & 57\,564 (69.6) & 39\,144 & 7\,947.76 & 6\,673.64 & 1\,551.10 & 9.53 \\
58 & 29 & 93\,870 & 77\,479 (82.5) & 30\,409 & 5\,272.34 & 5\,162.29 & 1\,434.49 & 7.23 \\
60 & 30 & 106\,504 & 74\,991 (70.4) & 48\,856 & 13\,191.39 & 9\,754.24 & 2\,418.80 & 14.67 \\
62 & 31 & 119\,984 & 99\,772 (83.2) & 37\,571 & 8\,719.97 & 8\,275.18 & 5\,020.18 & 13.09 \\
64 & 32 & 135\,116 & 96\,115 (71.1) & 60\,407 & 17\,674.78 & 13\,420.81 & 3\,270.29 & 22.58 \\
\midrule
Total & & 728\,182 & 552\,727 (75.9) & 290\,542 & 63\,657.71 & 55\,812.87 & 16\,214.97 & 80.34 \\
\bottomrule
\end{tabular}
\end{table}

The algebraic rules in Section~\ref{sec:algebraic} excluded
$552\,727$ of the $728\,182$ canonical cases (75.9\%).
Table~\ref{tab:pruning} gives the count by rule. The remaining cases
can each give more than one formula because of the branches in
Section~\ref{sec:branches}.

\begin{table}[!htbp]
\centering\small
\caption{Canonical cases excluded by each algebraic rule. The
  generator counts each excluded case under the first rule it fails.}
\label{tab:pruning}
\setlength{\tabcolsep}{10pt}
\renewcommand{\arraystretch}{1.12}
\begin{tabular}{@{}rrrrr@{}}
\toprule
$n$ & NS-1 & NS-2 & NS-3 & SQ-RC-Odd \\
\midrule
50 & 15\,064 & 24\,674 & 3\,731 & 819 \\
52 & 43\,354 & 0 & 0 & 0 \\
54 & 20\,265 & 32\,879 & 5\,005 & 1\,015 \\
56 & 57\,564 & 0 & 0 & 0 \\
58 & 26\,684 & 42\,975 & 6\,580 & 1\,240 \\
60 & 74\,991 & 0 & 0 & 0 \\
62 & 34\,544 & 55\,232 & 8\,500 & 1\,496 \\
64 & 96\,115 & 0 & 0 & 0 \\
\bottomrule
\end{tabular}
\end{table}

Each log records the formula tag and digest, the solver verdict,
the proof byte count, and the checker verdict. The Python auditor
\nolinkurl{scripts/verify_certificate.py} separately enumerates
the canonical cases, counts by pruning rule, and expected branch
tags. It checks that each expected tag appears exactly once,
reconciles the summary counts with the records, checks the logged
verdicts, and can verify upper-bound witnesses from their coordinates.
All $290\,542$ digests match regenerated CNFs, including the
$74\,155$ archived digests from the earlier $n=50,52,54$ runs.

The raw proofs, totaling $80.34$\,GB, were deleted after checking.
The driver keeps temporary files on verification failure or an
unexpected satisfiable result. The archived JSONL files are audit
records. To verify the proofs afresh, regenerate them from the
pinned code in Section~\ref{sec:data}.

\subsection{Timing and resource use}

For the eight runs in Table~\ref{tab:certified}, ten workers on
one Apple~M4 used CaDiCaL 3.1.0 (commit \texttt{e4d9bd93}).
The checker \texttt{drat-trim} (commit \texttt{2e3b2dc0}) reads
binary proofs with options \mbox{\texttt{-w -t 1000000}}.
The driver's \texttt{auto} mode selects CaDiCaL's \texttt{unsat}
preset for branch kinds 1, 2, 5, 7, 8, 11, and 12, and the default
preset otherwise.

The driver uses \texttt{std::chrono::steady\_clock}.
The solve column times the \texttt{solve()} calls, including DRAT output
during each call. The check column times checker invocations
from startup to exit. Wall time runs from the start of the
worker phase through encoding, file operations, solving, and
checking. Case enumeration and ordering precede that phase.
Solve and check times accumulate across workers and can therefore
exceed wall time.

For $n=56,58,60,62,64$, the script
\nolinkurl{scripts/run_certified_board.py} ran the boards in order.
It sampled free storage, resident set size (RSS), swap usage,
macOS VM pressure, and temporary file size every $0.5$\,s.
The monitor was set to suspend cleanly below $8$\,GiB of free
storage or at critical VM pressure. All five boards finished
without suspension, and free storage stayed above $21$\,GiB.

\begin{table}[!htbp]
\centering\small
\caption{Largest proof per board and recorded resource peaks.
  RSS sums the parent driver and its active checker processes.
  Temporary disk holds the active CNF and DRAT files.
  Dashes indicate runs without resource monitoring. The leaf
  tags omit the common frame prefix $(0,1,0,1)$ and suffix
  \texttt{\_0\_0\_0\_0\_0}.}
\label{tab:resources}
\setlength{\tabcolsep}{4.0pt}
\renewcommand{\arraystretch}{1.18}
\begin{tabular}{@{}rrrrlrrr@{}}
\toprule
$n$ & Peak RSS & Peak temp & Min free & Largest proof leaf & Leaf solve & Leaf check & Leaf DRAT \\
    & (MiB) & (MB) & disk (GiB) & (branch on $(0,1,0,1)$) & (s) & (s) & (MB) \\
\midrule
50 & --- & --- & --- & \texttt{0\_0\_0\_1\_b8} & 33.96 & 33.66 & 33.65 \\
52 & --- & --- & --- & \texttt{0\_0\_1\_1\_b3} & 23.12 & 21.57 & 29.83 \\
54 & --- & --- & --- & \texttt{0\_0\_0\_1\_b8} & 155.18 & 234.21 & 133.25 \\
\midrule
56 & 929.2 & 208.0 & 21.59 & \texttt{0\_0\_1\_1\_b3} & 32.59 & 40.14 & 73.53 \\
58 & 1\,769.5 & 683.0 & 21.16 & \texttt{0\_0\_0\_1\_b8} & 937.76 & 496.59 & 671.15 \\
60 & 1\,335.3 & 519.2 & 21.07 & \texttt{0\_0\_1\_1\_b3} & 267.03 & 326.73 & 369.37 \\
62 & 4\,896.0 & 3\,066.6 & 21.06 & \texttt{0\_0\_0\_1\_b8} & 2\,475.08 & 2\,544.16 & 3\,066.20 \\
64 & 3\,407.8 & 1\,953.9 & 22.04 & \texttt{0\_0\_1\_1\_b3} & 1\,185.09 & 1\,753.70 & 1\,830.66 \\
\bottomrule
\end{tabular}
\end{table}

The largest proofs all came from the frame $(0,1,0,1)$.
For $n\equiv2\pmod4$, the largest leaf was the one-diagonal-deficit
monochromatic branch \texttt{b8}. On $Q_{62}$ its proof took about
$5\,019$\,s to produce and check, almost all of the $5\,020.18$\,s
worker wall time. The other $37\,570$ leaves used a combined
$6\,244.89$\,s of solve time and $5\,731.02$\,s of check time
across workers. The largest proof occupied $3\,066.20$\,MB,
close to the $3\,066.6$\,MB peak for temporary files, and
peak RSS was $4\,896$\,MiB.

For $n\equiv0\pmod4$, the largest leaf was the two-diagonal-deficit
branch \texttt{b3}, which includes placements with either one or
two colors. Its proof grew from $29.83$\,MB at $n=52$ to
$1\,830.66$\,MB at $n=64$. The zero-deficit monochromatic leaf
\texttt{b12}, which uses the quadratic moment identity, took
less than $0.75$\,s to solve and at most $1.14$\,s including
checking on each tested board.

The worker queue gives priority to the $(0,1,0,1)$ leaves in
branches 8, 3, and 12. This schedules expensive proofs alongside
the bulk of the queue. The large \texttt{b3} leaf finished before
the end of the worker phase on both $Q_{60}$ and $Q_{64}$.

\subsection{Same-machine benchmarks on \texorpdfstring{$Q_{26}$}{Q26}}

To measure certification overhead on one Apple~M4 worker, we
selected up to 12 leaves per branch type in order of SHA-256
tag digest, for a total of 99 leaves. We used the pinned tools
and the \texttt{auto} configuration. We ran each mode three
times and alternated which mode went first in each pair. Median
worker wall time was $0.33$\,s without proofs and $3.35$\,s
with proof generation and checking. The sum of solver-call
times had medians of $0.18$\,s and $0.19$\,s, respectively.
In proof mode the sum of checker times had a median of
$2.88$\,s, including one invocation per leaf. All 99 proofs
passed in each of the three runs.

To measure the effect of algebraic pruning, we timed a separate
benchmark on the full set of $Q_{26}$ cases.
Both modes ran on one Apple~M4 with ten workers, the pinned
solver and checker, and identical solver and symmetry settings.
Only the algebraic pruning pass was toggled. The wall-clock
measurement covered the complete command, from case enumeration
and CNF construction through proof generation and checking.
Across three runs per mode in alternating order, pruning reduced
the formula count from $7\,668$ to $2\,773$ and median wall time
from $58.49$\,s to $24.39$\,s, a $2.40\times$ speedup on $Q_{26}$.
All proofs passed in all six runs.

\subsection{Structural results for larger boards}

The DRAT checks reported here use the encoding and coverage
proof in Sections~\ref{sec:prelim}--\ref{sec:encoding}. For
larger boards one could use Weakley's results to exclude whole
color classes before encoding them
\cite[Proposition~11, Theorems~17--18 and Section~3]{Weakley2022}.
For $n\equiv0\pmod4$, the only board sizes permitting a
monochromatic dominating $n/2$-set are $4$ and $12$. For
$n=2k>2$ with $k$ odd, use Weakley's standard orientation, with
coordinates $1-k,\ldots,k$ and all occupied rows and columns of
the same parity as $k$. In this orientation there are no excess
difference-diagonal indices and there are two excess sum-diagonal
indices, counted with multiplicity.

Let $d_{\mathrm W}$ be half the absolute difference of the two
excess sum indices and $e_{\mathrm W}=|s_0-s_2|/2$, where $s_r$ is the
least absolute value of a missing sum index in residue class $r$
modulo four. Weakley's necessary condition is
\[
  d_{\mathrm W}^{\,2}+(k-1)e_{\mathrm W}^{\,2}=\frac{k(k^2+2)}{3},
  \qquad d_{\mathrm W},e_{\mathrm W}\in\{1,3,\dots,k\}.
\]
For $2<n<238$, the equation leaves $n=6,10,18,30,38,54,130,142,234$.
Proposition~19 excludes $38,142$ ($d_{\mathrm W}=e_{\mathrm W}$),
and the negative $F(k)$ bounds in Weakley's Section~3 exclude
$54,234$. Taken together, these criteria rule out monochromatic
dominating $n/2$-sets on every even board from $50$ through $128$.
On those boards an encoder using the criteria could omit
monochromatic leaves such as \texttt{b8} and require both colors
in \texttt{b3}.

For bichromatic sets with $k>2$, Weakley~\cite[Proposition~11]{Weakley2022}
gives the count pair $\{z_0,z_1\}=\{\lfloor k/2\rfloor,\lceil k/2\rceil\}$,
with the additional pair $\{k/2-1,k/2+1\}$ possible only if $k,w,h$
are even and $d_R=d_C=0$. This would restrict the largest \texttt{b3}
leaf on $Q_{64}$ to $z_0\in\{15,16,17\}$. For nonsquare
even-by-even frames with $w>h$, one can intersect those count bounds
with $z_{\bar q_*}\le h/2$. The margin sum in Lemma~\ref{lem:ns1}(ii)
and $w\le2k-2$ already give $h\ge k$.

\section{Code and data availability}\label{sec:data}

The encoder, checker wrapper, runner, logs, diagnostics, and
witness coordinates are available at
\href{https://github.com/CosmicCrusader23/queens}{\nolinkurl{https://github.com/CosmicCrusader23/queens}}.
The pinned
\href{https://github.com/CosmicCrusader23/queens/blob/\ReproCommit/README.md}{README}
gives the full build, audit, regeneration, and benchmark commands.

\begin{itemize}[leftmargin=1.8em,itemsep=2pt]
  \item Reproduction commit
    \href{https://github.com/CosmicCrusader23/queens/tree/\ReproCommit}{\nolinkurl{61e8df4438cad8837f4498ad4160394028ca7232}}.
  \item Earlier $Q_{50},Q_{52},Q_{54}$ snapshot
    \href{https://github.com/CosmicCrusader23/queens/tree/7872bd98a9d1d72d463c6b967717bd7ab093b7d8}{\nolinkurl{7872bd98a9d1d72d463c6b967717bd7ab093b7d8}}.
  \item CaDiCaL~3.1.0 commit
    \href{https://github.com/arminbiere/cadical/commit/e4d9bd93af444c6575a908a24da2c4130cb2a1c7}{\nolinkurl{e4d9bd93af444c6575a908a24da2c4130cb2a1c7}},
  \item \texttt{drat-trim} commit
    \href{https://github.com/marijnheule/drat-trim/commit/2e3b2dc0ecf938addbd779d42877b6ed69d9a985}{\nolinkurl{2e3b2dc0ecf938addbd779d42877b6ed69d9a985}}.
\end{itemize}
The runs used a ten-core Apple~M4 (arm64, 16\,GB RAM) running
macOS~26.5.2 (Darwin~25.5.0) with Apple Clang~17.0.0, C++20, and
Python~3.14.0. The Python scripts use only the standard library.

After the dependency setup and \texttt{make -j4}, the following
commands regenerate and check the $Q_{56}$ proofs in a fresh
output directory.
\begin{quote}\small
\begin{verbatim}
mkdir -p regeneration
python3 scripts/run_certified_board.py --n 56 --jobs 10 \
  --out regeneration/cert_n56_k28.jsonl \
  --diag-out regeneration/diagnostics_n56_k28.json
\end{verbatim}
\end{quote}
For other sizes in Table~\ref{tab:certified}, change \texttt{-{}-n}
and the output filenames. A direct driver invocation uses the
same board size and \texttt{-{}-jobs 10 -{}-certify -{}-config auto}.
Table~\ref{tab:certified} gives the expected leaf counts and proof
volumes. If retaining proofs with \texttt{-{}-proof-dir}, budget
for the full per-board volume. The monitor above runs with
temporary proof files.

For an audit of the archived logs, use the driver's
\texttt{-{}-fingerprints} option and pass the output and the
matching JSONL log to \nolinkurl{scripts/verify_certificate.py}
with \texttt{-{}-fingerprints} and \texttt{-{}-cert}, respectively,
along with the board size in \texttt{-{}-n}.
The same auditor can check the placements directly with
\texttt{-{}-witness-json} \nolinkurl{results/even_upper_witnesses.json}.

\section{AI contributions}\label{sec:provenance}

Gemini~4~Neon, used through Google Antigravity, developed the
mathematical reduction, wrote the C++ and Python code, ran the
computations and checks, and drafted the manuscript. OpenAI
Codex reviewed the proofs and encoder, checked reproduction of
the computations, and helped revise the manuscript. The author
initiated and coordinated the project and reviewed the
computation logs and verification outputs.

\appendix
\section{Upper-bound placements}\label{sec:witnesses}

Write $(x,y)=(\text{column},\text{row})$, starting at zero.
The search in \nolinkurl{src/witness_sa.cpp} placed queens on a
permutation of the odd rows and columns plus one extra square.
The verifier checks the size and coverage of every placement
below. It also checks distinct rows, columns, and diagonals
for the independent placements.

\subsection{Even boards from 50 through 64}

These placements give the upper bounds in Theorem~\ref{thm:main}.
For Corollary~\ref{cor:indep} we separately checked the eight
independent placements of \"Osterg\aa rd and Weakley
\cite[Section~5, p.~11]{OW2001} in their original coordinates.
For $Q_{54}$, their listed queens dominate $[-27,26]\times[-26,27]$,
which adjoins column $-27$ and row $27$ to $Q_{53}$.
Their $Q_{58}$ construction adjoins row and column $-29$ to $Q_{57}$
and adds $(17,-29)$, giving the board $[-29,28]^2$.

\begin{placement}{$Q_{50}$ ($26$ queens)}
$(0,0)$, $(1,3)$, $(3,35)$, $(5,25)$, $(7,15)$, $(9,9)$, $(11,39)$, $(13,49)$, $(15,19)$, $(17,27)$, $(19,11)$, $(21,45)$, $(23,29)$, $(25,1)$, $(27,21)$, $(29,41)$, $(31,47)$, $(33,23)$, $(35,7)$, $(37,17)$, $(39,37)$, $(41,5)$, $(43,31)$, $(45,13)$, $(47,43)$, $(49,33)$.
\end{placement}

\begin{placement}{$Q_{52}$ ($27$ queens)}
$(0,0)$, $(1,15)$, $(3,29)$, $(5,43)$, $(7,37)$, $(9,41)$, $(11,45)$, $(13,11)$, $(15,7)$, $(17,19)$, $(19,1)$, $(21,25)$, $(23,31)$, $(25,3)$, $(27,49)$, $(29,39)$, $(31,27)$, $(33,51)$, $(35,5)$, $(37,23)$, $(39,33)$, $(41,47)$, $(43,9)$, $(45,35)$, $(47,21)$, $(49,17)$, $(51,13)$.
\end{placement}

\begin{placement}{$Q_{54}$ ($28$ queens)}
$(0,0)$, $(1,13)$, $(3,53)$, $(5,25)$, $(7,15)$, $(9,37)$, $(11,23)$, $(13,45)$, $(15,51)$, $(17,1)$, $(19,7)$, $(21,17)$, $(23,47)$, $(25,27)$, $(27,21)$, $(29,35)$, $(31,29)$, $(33,49)$, $(35,39)$, $(37,5)$, $(39,11)$, $(41,3)$, $(43,19)$, $(45,9)$, $(47,31)$, $(49,41)$, $(51,43)$, $(53,33)$.
\end{placement}

\begin{placement}{$Q_{56}$ ($29$ queens)}
$(0,0)$, $(1,17)$, $(3,23)$, $(5,45)$, $(7,51)$, $(9,33)$, $(11,43)$, $(13,19)$, $(15,7)$, $(17,13)$, $(19,55)$, $(21,31)$, $(23,21)$, $(25,53)$, $(27,29)$, $(29,5)$, $(31,39)$, $(33,27)$, $(35,47)$, $(37,1)$, $(39,11)$, $(41,25)$, $(43,3)$, $(45,49)$, $(47,15)$, $(49,37)$, $(51,41)$, $(53,9)$, $(55,35)$.
\end{placement}

\begin{placement}{$Q_{58}$ ($30$ queens)}
$(0,0)$, $(1,35)$, $(3,37)$, $(5,15)$, $(7,21)$, $(9,39)$, $(11,5)$, $(13,57)$, $(15,53)$, $(17,7)$, $(19,27)$, $(21,23)$, $(23,41)$, $(25,51)$, $(27,31)$, $(29,25)$, $(31,1)$, $(33,55)$, $(35,45)$, $(37,29)$, $(39,13)$, $(41,19)$, $(43,49)$, $(45,11)$, $(47,3)$, $(49,47)$, $(51,33)$, $(53,9)$, $(55,17)$, $(57,43)$.
\end{placement}

\begin{placement}{$Q_{60}$ ($31$ queens, independent)}
$(0,0)$, $(1,45)$, $(3,39)$, $(5,15)$, $(7,27)$, $(9,13)$, $(11,43)$, $(13,3)$, $(15,11)$, $(17,57)$, $(19,47)$, $(21,9)$, $(23,37)$, $(25,31)$, $(27,25)$, $(29,53)$, $(31,17)$, $(33,35)$, $(35,29)$, $(37,1)$, $(39,55)$, $(41,49)$, $(43,19)$, $(45,5)$, $(47,59)$, $(49,21)$, $(51,7)$, $(53,33)$, $(55,23)$, $(57,41)$, $(59,51)$.
\end{placement}

\begin{placement}{$Q_{62}$ ($32$ queens)}
$(0,0)$, $(1,15)$, $(3,21)$, $(5,47)$, $(7,33)$, $(9,19)$, $(11,9)$, $(13,43)$, $(15,55)$, $(17,45)$, $(19,57)$, $(21,39)$, $(23,27)$, $(25,59)$, $(27,5)$, $(29,25)$, $(31,17)$, $(33,41)$, $(35,1)$, $(37,29)$, $(39,61)$, $(41,3)$, $(43,37)$, $(45,51)$, $(47,11)$, $(49,23)$, $(51,53)$, $(53,35)$, $(55,13)$, $(57,7)$, $(59,49)$, $(61,31)$.
\end{placement}

\begin{placement}{$Q_{64}$ ($33$ queens)}
$(0,0)$, $(1,41)$, $(3,23)$, $(5,53)$, $(7,43)$, $(9,25)$, $(11,7)$, $(13,17)$, $(15,59)$, $(17,5)$, $(19,11)$, $(21,31)$, $(23,37)$, $(25,57)$, $(27,51)$, $(29,9)$, $(31,33)$, $(33,61)$, $(35,21)$, $(37,35)$, $(39,29)$, $(41,49)$, $(43,55)$, $(45,1)$, $(47,19)$, $(49,13)$, $(51,3)$, $(53,47)$, $(55,15)$, $(57,63)$, $(59,27)$, $(61,45)$, $(63,39)$.
\end{placement}

\subsection{Earlier independent values}

\begin{theorem}[{Neuhaus~\cite[Section~2.4]{Neuhaus2009}}]\label{thm:indep}
The independent domination numbers satisfy
\[
  i(Q_{32})=17,\quad i(Q_{34})=18,\quad i(Q_{38})=20,\quad
  i(Q_{40})=21,\quad i(Q_{44})=23.
\]
\end{theorem}

\begin{proof}
The matching ordinary domination numbers come from
\emph{Hint and Conquer}~\cite{Hint2026} for $n=32,34,38$ and
Davies~\cite[Theorem~1.3 and Table~2]{Davies2026} for $n=40,44$.
The independent placements below attain these lower bounds.
\end{proof}

\begin{placement}{$Q_{32}$ ($17$ queens)}
$(1,25)$, $(2,2)$, $(3,19)$, $(5,9)$, $(7,27)$, $(9,1)$, $(11,17)$, $(13,11)$, $(15,3)$, $(17,29)$, $(19,21)$, $(21,15)$, $(23,31)$, $(25,5)$, $(27,23)$, $(29,13)$, $(31,7)$.
\end{placement}

\begin{placement}{$Q_{34}$ ($18$ queens)}
$(0,32)$, $(1,9)$, $(3,19)$, $(5,25)$, $(7,31)$, $(9,5)$, $(11,17)$, $(13,11)$, $(15,3)$, $(17,29)$, $(19,21)$, $(21,15)$, $(23,27)$, $(25,1)$, $(27,7)$, $(29,13)$, $(31,23)$, $(33,33)$.
\end{placement}

\begin{placement}{$Q_{38}$ ($20$ queens)}
$(1,19)$, $(3,11)$, $(4,34)$, $(5,27)$, $(7,5)$, $(9,35)$, $(11,13)$, $(13,17)$, $(15,1)$, $(17,31)$, $(19,15)$, $(21,21)$, $(23,33)$, $(25,3)$, $(27,9)$, $(29,23)$, $(31,37)$, $(33,7)$, $(35,25)$, $(37,29)$.
\end{placement}

\begin{placement}{$Q_{40}$ ($21$ queens)}
$(0,38)$, $(1,35)$, $(3,19)$, $(5,13)$, $(7,7)$, $(9,1)$, $(11,23)$, $(13,37)$, $(15,29)$, $(17,15)$, $(19,39)$, $(21,5)$, $(23,17)$, $(25,27)$, $(27,3)$, $(29,33)$, $(31,11)$, $(33,21)$, $(35,31)$, $(37,9)$, $(39,25)$.
\end{placement}

\begin{placement}{$Q_{44}$ ($23$ queens)}
$(0,42)$, $(1,37)$, $(3,13)$, $(5,7)$, $(7,25)$, $(9,31)$, $(11,41)$, $(13,11)$, $(15,19)$, $(17,43)$, $(19,1)$, $(21,35)$, $(23,23)$, $(25,3)$, $(27,17)$, $(29,21)$, $(31,5)$, $(33,39)$, $(35,29)$, $(37,33)$, $(39,9)$, $(41,27)$, $(43,15)$.
\end{placement}

\clearpage
\begin{thebibliography}{99}
\small
\raggedright
\setlength{\itemsep}{4pt plus 1pt}

\bibitem{BB2003}
O.~Bailleux and Y.~Boufkhad,
\emph{Efficient CNF encoding of Boolean cardinality constraints},
in: Principles and Practice of Constraint Programming (CP 2003),
Lecture Notes in Comput.\ Sci.\ \textbf{2833}, Springer, 2003,
pp.~108--122.
\href{https://doi.org/10.1007/978-3-540-45193-8_8}{DOI}.

\bibitem{CaDiCaL}
A.~Biere, T.~Faller, K.~Fazekas, M.~Fleury, N.~Froleyks, and F.~Pollitt,
\emph{CaDiCaL~2.0},
in: Computer Aided Verification (CAV 2024),
Lecture Notes in Comput.\ Sci.\ \textbf{14681}, Springer, 2024,
pp.~133--152.
\href{https://doi.org/10.1007/978-3-031-65627-9_7}{DOI}.

\bibitem{Cockayne1990}
E.~J.~Cockayne,
\emph{Chessboard domination problems},
Discrete Math.\ \textbf{86} (1990), 13--20.
\href{https://doi.org/10.1016/0012-365X(90)90344-H}{DOI}.

\bibitem{Davies2026}
J.~Davies,
\emph{$\gamma(Q_{35})=19$ and certified queen domination numbers up to $n=46$},
Hexagon Math, 2026.
\href{https://hexagonmath.org/2611.00830v1}{2611.00830v1}.

\bibitem{Djukic2026}
Aleksej \DJ uki\'c,
comment on $\gamma(Q_{48})$ in OEIS sequence A075458,
November 22, 2026.
\href{https://oeis.org/A075458}{Sequence}
and \href{https://oeis.org/A075458/a075458_1.txt}{computation summary and witness}.

\bibitem{Guy1994}
R.~K.~Guy,
\emph{Unsolved Problems in Number Theory},
2nd ed., Springer, New York, 1994, Problem~C18.
\href{https://doi.org/10.1007/978-1-4899-3585-4}{DOI}.

\bibitem{BSC2026}
jkolantree,
\emph{$Q_{26}$ grid-annihilator formalization},
Lean development in the \texttt{BSC} repository, 2026.
\href{https://github.com/jkolantree/BSC/tree/main/formal/q26_grid_annihilator}{Project and theorem statement}.

\bibitem{Neuhaus2009}
S.~Neuhaus,
\emph{Eigenschaften kleinster dominierender Mengen und Dominanzzahlen von Damengraphen},
Ph.D.\ dissertation, Universit\"at zu K\"oln, 2009.
\href{https://kups.ub.uni-koeln.de/2916/}{University repository}.

\bibitem{OW2001}
P.~R.~J.~\"Osterg\aa rd and W.~D.~Weakley,
\emph{Values of domination numbers of the queen's graph},
Electron.\ J.\ Combin.\ \textbf{8}(1) (2001), \#R29.
\href{https://doi.org/10.37236/1573}{DOI}.

\bibitem{RV1987}
V.~Raghavan and S.~M.~Venkatesan,
\emph{On bounds for a board covering problem},
Inform.\ Process.\ Lett.\ \textbf{25}(5) (1987), 281--284.
\href{https://doi.org/10.1016/0020-0190(87)90201-8}{DOI}.

\bibitem{Hint2026}
T.~Rostami,
\emph{Hint and Conquer: New Queen Domination Numbers by SAT Solving},
preprint, 2026.
\href{https://github.com/TahaRostami/Gamma/blob/hint-and-conquer/doc.pdf}{Manuscript}.
Author attribution follows \href{https://oeis.org/A075458}{OEIS~A075458}.

\bibitem{Sinz2005}
C.~Sinz,
\emph{Towards an optimal CNF encoding of Boolean cardinality constraints},
in: Principles and Practice of Constraint Programming (CP 2005),
Lecture Notes in Comput.\ Sci.\ \textbf{3709}, Springer, 2005,
pp.~827--831.
\href{https://doi.org/10.1007/11564751_73}{DOI}.

\bibitem{Weakley2022}
W.~D.~Weakley,
\emph{Queen domination of even square boards},
Electron.\ J.\ Combin.\ \textbf{29}(2) (2022), \#P2.50.
\href{https://doi.org/10.37236/10617}{DOI}.

\bibitem{DRATtrim}
N.~Wetzler, M.~J.~H.~Heule, and W.~A.~Hunt~Jr.,
\emph{DRAT-trim: Efficient checking and trimming using expressive clausal proofs},
in: Theory and Applications of Satisfiability Testing (SAT 2014),
Lecture Notes in Comput.\ Sci.\ \textbf{8561}, Springer, 2014,
pp.~422--429.
\href{https://doi.org/10.1007/978-3-319-09284-3_31}{DOI}.

\end{thebibliography}

\end{document}
